% =====================================================================
%  18.4  Fenwick on Trees: driving Ch.5's Euler array   sec:c14-trees
%  Exports: def:c14-eulerbit  prop:c14-eulerrange  prop:c14-path
%           thm:c14-parity    lst:c14-flatten  lst:c14-subtree
%           lst:c14-path  lst:c14-383-soln  fig:c14-euler  prob:c14-prop
%  Reuses (never re-derives): Ch.5 (\chref{c24}) tin/tout + subtree=interval
%    (\thmref{c24-range}); \secref{c14-core} plain BIT; \secref{c14-range}
%    difference BIT (range add / point read).
%  Running tree T (Ch.12--17), root 1, children ascending, edges
%    1-2,1-3,2-4,2-5,3-6,3-7,5-8,7-9,9-10.  Ch.16 preorder (INCLUSIVE tout):
%      tin  (by node 1..10) = (1,2,6,3,4,7,8,5,9,10)
%      tout (by node 1..10) = (10,5,10,3,5,7,10,5,10,10)
%      size=(10,4,5,1,2,1,3,1,2,1)  depth=(0,1,1,2,2,2,2,3,3,4)
%      Euler order (position -> node) = (1,2,4,5,8,3,6,7,9,10)
%      sgn (depth parity, by node)    = (+1,-1,-1,+1,+1,+1,+1,-1,-1,+1)
%  Opener six-angle row (BIT on T) reproduced in the subtree-add trace:
%    subtree(2)=[2,5] add +3, then whole tree +2: node4 reads 5, node6 reads 2.
%  CF 383C Propagating tree (2000) -- live-verified via the CF API + statement:
%    n,m<=200000, a_i<=1000, val<=1000, root 1.
%    Query "1 x val": add val to x, then -val to children, +val to grand-
%    children, ... alternating by depth through the whole subtree.
%    Query "2 x": print current value at x.
%    Sample (verified): 5 5 / 1 2 1 1 2 / edges 1-2,1-3,2-4,2-5 /
%      1 2 3 / 1 1 2 / 2 1 / 2 2 / 2 4  => output 3 / 3 / 0.
%    Solution: single difference BIT over Euler positions; store delta=val*sgn[x]
%      on range [tin[x],tout[x]]; read a[x] + sgn[x]*point(tin[x]).
% =====================================================================
\section{Fenwick on Trees: Driving the Euler Array}\label{sec:c14-trees}

\begin{previouslybox}[Ch.~5 flattened the tree; it just had no engine yet]
Chapter~5 (\chref{c24}) proved the reduction we now exploit without re-deriving
it: a depth-first \term{Euler tour} stamps each node \(v\) an entry time
\(\tin(v)\) and an exit time \(\tout(v)\) so that \emph{the subtree of \(v\) is
exactly the index interval} \([\tin(v),\tout(v)]\) (\thmref{c24-range}).  What
Ch.~5 could do was read static sums over that interval with a prefix array
(\(\propref{c24-prefix}\)); what it \emph{could not} do was keep answering
after the values changed, because a prefix array dies under updates.  This
section bolts \secref{c14-core}'s BIT onto Ch.~5's flat array, discharging the
forward pointer Ch.~5 left here: subtree ranges now get an engine.
\end{previouslybox}

\begin{topicintro}
Once the tree is flattened, every ``add to a subtree / read a node / sum a
subtree'' question stops being a \emph{tree} problem and becomes a \emph{range}
problem over the slots \(1..n\) -- and ranges over slots with an invertible
accumulator are exactly what a BIT is built for.  Concretely we place each
node's value at slot \(\tin(v)\), then reach for whichever \secref{c14-range}
flavour the query shape forces: a \emph{plain} BIT for point-update + subtree
sum, a \emph{difference} BIT for subtree-add + node-read, and the
\emph{enter\(+1\)/leave\(-1\)} re-interpretation of the difference BIT for
root-path reads.  The single new idea -- used to crack the solved problem -- is
that a sign that alternates with \emph{depth} becomes a \(\pm1\) multiplier
attached to each slot, so a depth-propagating update is still \emph{one} range
addition.

\textbf{Prerequisites.} Ch.~5 (\chref{c24}) for \(\tin/\tout\);
\secref{c14-core} for the two walks; \secref{c14-range} for the difference BIT
(range add / point read).  The running example is the tree \(T\) with the
stamps recorded in \figref{c14-array}'s opener block.

\textbf{After this section you can}
\begin{itemize}[nosep]
  \item turn any subtree update/query into a range operation on the Euler
        array and pick the right BIT flavour; \hfill\lv{3}
  \item answer root-path reads with the two-event \(\mathit{enter}{+}1 /
        \mathit{leave}{-}1\) trick, no LCA machinery; \hfill\lv{4}
  \item absorb a depth-parity (or any per-node \(\pm1\)) into the stored
        value so a propagating update stays a single range add; \hfill\lv{4}--\lv{5}
  \item solve CF~383C with one difference BIT over the tour. \hfill\lv{5}
\end{itemize}
\end{topicintro}

\begin{recognitionbox}[Pattern card -- ``many subtree / root-path operations on one static tree'']
\reckey{``add \(x\) to every node in the subtree of \(v\)'' interleaved with
``read node \(v\)'`; ``invert the colour of a whole subtree, query one
vertex''; ``sum of the values on the path from the root to \(v\)''; ``the
value at \(v\) after updates that propagate to descendants with alternating
sign''; ``how many marked nodes are ancestors of \(v\)''.  The tree itself
never changes -- only per-node numbers do, and queries cluster on
subtrees/paths.}
\recuse{
  \item subtree add + point read (or point add + subtree sum) -- flatten to
        \([\tin,\tout]\), difference BIT (this section, CF~383C, CSES~1137);
  \item root-path read with point ``activation'' -- enter\(+1\)/leave\(-1\)
        on the same difference BIT (\(\propref{c14-path}\));
  \item updates whose effect depends on depth parity or a per-node sign --
        fold the sign into the stored delta (\(\thmref{c14-parity}\)).
}
\recnotuse{
  \item the tree is \emph{cut, linked or rerooted} between queries --
        Euler-tour trees / splay, Ch.~6 (\chref{c06}); a static tour array
        cannot survive a structural change;
  \item the range aggregate is not invertible (max/min/gcd over the subtree)
        -- segment tree with the flat array, Ch.~19 (\chref{c13});
  \item paths are \emph{arbitrary} \(u\)-to-\(v\) with non-commutative data
        -- heavy-light decomposition on top of a segment tree, not a BIT;
  \item the tree is a single path -- it is just an array; use
        \secref{c14-range} directly.
}
One-line test: \emph{is the tree static and does every operation address a
subtree or a root-path?  Flatten with Ch.~5's \(\tin/\tout\), then let a BIT
own the range.}
\end{recognitionbox}

% ---------------------------------------------------------------
\subsection{The Flattened Array and Its Two Range Shapes}
% ---------------------------------------------------------------

\begin{definition}{Euler-indexed BIT}{c14-eulerbit}
Fix a DFS order of the tree and let \(\mathit{euler}[p]\) be the node stamped
\(\tin(\mathit{euler}[p])=p\).  An \term{Euler-indexed BIT} is a BIT
\(B[1..n]\) whose slot \(p\) stores the aggregate of the \emph{node}
\(\mathit{euler}[p]\).  A node operation on \(v\) becomes a slot operation at
\(\tin(v)\); a subtree operation on \(v\) becomes a range operation on
\([\tin(v),\tout(v)]\) (both \(1\)-based, \(\tout\) \emph{inclusive}, per this
chapter's convention).  Because \(\tin\) is a permutation, ``node \(v\)'' and
``slot \(\tin(v)\)'' name the same object -- we translate silently between
them.
\end{definition}

\begin{figure}[H]
\centering
\begin{tikzpicture}[diagram, scale=0.62, transform shape,
    tnode/.style={circle, draw=primaryblue, line width=0.8pt, fill=primaryblue!8,
                  minimum size=6.5mm, font=\scriptsize\bfseries, text=primaryblue!55!black}]
  % tree edges (root 1 at top), children ascending left-to-right
  \node[tnode] (n1) at (0,3)   {1};
  \node[tnode] (n2) at (-3,2)  {2};
  \node[tnode] (n3) at (3,2)   {3};
  \node[tnode] (n4) at (-4,1)  {4};
  \node[tnode] (n5) at (-1.5,1){5};
  \node[tnode] (n6) at (2,1)   {6};
  \node[tnode] (n7) at (4.2,1) {7};
  \node[tnode] (n8) at (-1.5,0){8};
  \node[tnode] (n9) at (4.2,0) {9};
  \node[tnode] (n10) at (4.2,-1){10};
  \draw (n1)--(n2) (n1)--(n3) (n2)--(n4) (n2)--(n5) (n3)--(n6) (n3)--(n7)
        (n5)--(n8) (n7)--(n9) (n9)--(n10);
  % tin/tout annotations
  \node[annot, font=\tiny, anchor=south] at (n1.north) {\(1[1,\!10]\)};
  \node[annot, font=\tiny, anchor=east]  at (n2.west)  {\(2[2,\!5]\)};
  \node[annot, font=\tiny, anchor=east]  at (n3.west)  {\(3[6,\!10]\)};
  \node[annot, font=\tiny, anchor=south] at (n4.north) {\(4[3]\)};
  \node[annot, font=\tiny, anchor=north] at (n5.south) {\(5[4,\!5]\)};
  \node[annot, font=\tiny, anchor=south] at (n6.north) {\(6[7]\)};
  \node[annot, font=\tiny, anchor=south] at (n7.north) {\(7[8,\!10]\)};
  \node[annot, font=\tiny, anchor=east]  at (n8.west)  {\(8[5]\)};
  \node[annot, font=\tiny, anchor=east]  at (n9.west)  {\(9[9,\!10]\)};
  \node[annot, font=\tiny, anchor=west]  at (n10.east) {\(10[10]\)};
  % Euler slot strip below the tree
  \foreach \p/\nd in {1/1,2/2,3/4,4/5,5/8,6/3,7/6,8/7,9/9,10/10} {
    \draw[difforange, line width=0.6pt, fill=difforange!8] (\p*0.62,-2.6) rectangle (\p*0.62+0.62,-2.0);
    \node[font=\tiny\ttfamily, text=orange!55!black] at (\p*0.62+0.31,-2.3) {\nd};
    \node[font=\tiny\ttfamily, text=darkgray] at (\p*0.62+0.31,-1.82) {\p};}
  \node[font=\scriptsize, text=orange!55!black, anchor=east] at (0.4,-2.3) {euler:};
  % subtree(2)=[2,5] and subtree(3)=[6,10] interval bars
  \fill[successgreen!35] (2*0.62+0.05,-3.0) rectangle (5*0.62+0.57,-2.75);
  \node[font=\tiny, text=green!45!black] at (3.5*0.62+0.3,-3.15) {subtree(2)=[2,5]};
  \fill[softpurple!28] (6*0.62+0.05,-3.0) rectangle (10*0.62+0.57,-2.75);
  \node[font=\tiny, text=purple!55!black] at (8*0.62+0.3,-3.15) {subtree(3)=[6,10]};
\end{tikzpicture}
\caption{Tree \(T\) with each node annotated \emph{v}[\(\tin(v),\tout(v)\)]
(Ch.~5's tour, \(\tout\) inclusive).  The orange strip is the Euler array:
slot~\(p\) holds node \(\mathit{euler}[p]\), read as \((1,2,4,5,8,3,6,7,9,10)\).
The two shaded bars are the intervals a BIT is handed to answer ``sum/add
subtree(2)'' and ``sum/add subtree(3)'' -- the whole tree collapses to
contiguous slots.  Node~3 sits at slot~6, so it is \emph{outside} the
subtree(2) bar even though it is drawn beside node~2 in the picture.}
\label{fig:c14-euler}
\end{figure}

Two operation shapes cover almost every tree-BIT problem.  Both are just a
flavour from \secref{c14-core}/\secref{c14-range} wearing the \(\tin\)
translation:

\begin{proposition}{Subtree sum with point updates; subtree add with point reads}{c14-eulerrange}
Let \(B\) be a BIT over the Euler array.
\begin{itemize}[nosep,leftmargin=1.4em]
  \item \emph{Point update of a node, subtree sum query} (plain BIT,
        \secref{c14-core}): to add \(\delta\) to node \(v\), run
        \(\code{update}(\tin(v),\delta)\); the sum of the whole subtree of
        \(u\) is \(\code{query}(\tout(u))-\code{query}(\tin(u)-1)\).
        This is CSES~1137 (``Subtree Queries''), solved in Ch.~5 with a static
        prefix array; here updates are free.
  \item \emph{Subtree add, single-node read} (difference BIT,
        \secref{c14-range}): to add \(\delta\) to every node of
        \(\subtree(v)\), run \(\code{range\_add}(\tin(v),\tout(v),\delta)\);
        the current value at node \(u\) is the point read
        \(\code{point}(\tin(u))\).  Cost \(\bigO{\log n}\) each.
\end{itemize}
\end{proposition}

\begin{lstlisting}[caption={Euler-indexed BIT: subtree sum with point updates, and subtree add with point reads}, label={lst:c14-subtree}]
// Ch.5 produced tin[1..n], tout[1..n]; we never recompute them here.
int lowbit(int i){ return i & -i; }

// (A) PLAIN BIT: add delta to node v, read the sum over subtree(u).
void bit_add(long long B[], int i, long long d){ for(; i<=n; i+=lowbit(i)) B[i]+=d; }
long long bit_pre(long long B[], int i){
    long long s=0; for(; i>0; i-=lowbit(i)) s+=B[i]; return s; }
void update_node(int v, long long d){ bit_add(B, tin[v], d); }          // 18.1
long long subtree_sum(int u){                                          // 18.1
    return bit_pre(B, tout[u]) - bit_pre(B, tin[u]-1); }

// (B) DIFFERENCE BIT: add d to every node of subtree(v), read one node.
void diff_add(long long D[], int i, long long d){ for(; i<=n; i+=lowbit(i)) D[i]+=d; }
void subtree_add(int v, long long d){        // range [tin, tout] = subtree(v)
    diff_add(D, tin[v], d); diff_add(D, tout[v]+1, -d); }   // 18.2 (i>n is a no-op)
long long node_value(int u){                 // point read at slot tin[u]
    long long s=0; for(int i=tin[u]; i>0; i-=lowbit(i)) s+=D[i]; return s; }
\end{lstlisting}

\begin{tracebox}[Subtree add on \(T\): \(\subtree(2){+}3\), then whole tree \(+2\)]
All nodes start at \(0\).  \(\subtree(2)=[\tin_2,\tout_2]=[2,5]\); the whole
tree is \(\subtree(1)=[1,10]\).  Two difference-BIT range adds:
\(\text{diff\_add}(2,+3),\text{diff\_add}(6,-3)\) then
\(\text{diff\_add}(1,+2),\text{diff\_add}(11,-2)\) (the last is \(i>n\), a
no-op).  Reading \(\code{node\_value}\) as a prefix sum over the difference
array at each node's \(\tin\):
\[
\renewcommand{\arraystretch}{1.1}
\begin{array}{c|c|c|c|c}
\text{node }u & \tin(u) & \tin\in[2,5]? & \tin\in[1,10]? & \text{value}\\ \hline
4 & 3 & \checkmark\ (+3) & \checkmark\ (+2) & \mathbf{5}\\
6 & 7 & \times & \checkmark\ (+2) & \mathbf{2}\\
2 & 2 & \checkmark\ (+3) & \checkmark\ (+2) & 5\\
3 & 6 & \times & \checkmark\ (+2) & 2
\end{array}
\]
Node~4 (a child of~2) carries both adds \(\to5\); node~6 (under~3) misses the
subtree(2) add \(\to2\).  This is exactly the opener's \textsc{BIT on \(T\)}
row -- and \figref{c14-euler}'s interval bars make it obvious without any
walking.
\end{tracebox}

\begin{complexitybox}[Why a BIT and not a segment tree here]
Both shapes of \(\propref{c14-eulerrange}\) are \(\bigO{\log n}\) with a
\emph{single} array and four-line code.  A segment tree does the same in the
same asymptotics but carries lazy tags and a \(4n\) array; the BIT is smaller,
faster, and cache-friendly.  The BIT is the right tool \emph{because} the
aggregate is a sum (invertible).  The moment the subtree query wants a
non-invertible combine -- ``maximum value in the subtree'' or ``first node
passing a threshold with arbitrary decrements'' -- the flat interval no longer
supports the difference trick and you owe the range to Ch.~19's segment tree
(\chref{c13}).  That boundary is stated here, not coded here.
\end{complexitybox}

% ---------------------------------------------------------------
\subsection{Root-Path Reads: the Enter\(	ext{+}1\) / Leave\(	ext{-}1\) Trick}
% ---------------------------------------------------------------

The difference BIT already stores a range add as a \(+\delta\) entering the
interval and a \(-\delta\) leaving it.  Read that backwards: if we \emph{tag}
a node \(v\) with a value \(x\) by writing \(+x\) at \(\tin(v)\) and \(-x\) at
\(\tout(v){+}1\), then a point read at \(\tin(u)\) sums \(x\) over exactly the
tags whose interval contains \(\tin(u)\) -- and \(\tin(v)\le\tin(u)\le\tout(v)\)
holds \emph{iff} \(v\) is an ancestor of \(u\) (Ch.~5, \(\thmref{c24-range}\)).
So the point read is the \emph{root-path} sum.

\begin{proposition}{Point-read at \(\tin(u)\) equals the root-to-\(u\) path sum}{c14-path}
Maintain a difference BIT \(D\).  To ``activate'' node \(v\) with weight \(x\)
for all its descendants, do \(D[\tin(v)]{+}{=}x,\;D[\tout(v){+}1]{-}{=}x\).
Then \(\sum_{p\le\tin(u)} D[p]\) (the point read at node \(u\)) equals
\(\sum_{v \text{ ancestor of } u} x_v\): the running prefix at slot
\(\tin(u)\) is inside \(v\)'s live window precisely when \(v\) is on the path
from the root to \(u\).  Each activation and each path read is \(\bigO{\log n}\).
\end{proposition}

\begin{lstlisting}[caption={two-event activation: subtree-add semantics become root-path reads}, label={lst:c14-path}]
// +x on ENTERING v's window, -x on LEAVING it -- one range add over [tin, tout].
void activate(int v, long long x){
    diff_add(D, tin[v],     x);        // switch x ON for v's whole subtree
    diff_add(D, tout[v] + 1, -x);      // switch it OFF after the subtree
}
// point read at tin[u] = sum of x over every ancestor v of u (root path of u):
long long root_path_sum(int u){
    long long s=0; for(int i=tin[u]; i>0; i-=lowbit(i)) s+=D[i]; return s; }
// So "add x to subtree(v); read value at u" and
//    "activate v with x; read the root-to-u sum" are the SAME two numbers
//    (x, -x) at (tin(v), tout(v)+1) -- just read with opposite intent.
\end{lstlisting}

\begin{insightbox}[One pair of writes, two stories]
The lines \(\mathit{diff\_add}(D,\tin(v),x)\) and
\(\mathit{diff\_add}(D,\tout(v){+}1,-x)\) are the same whether you call them
``add \(x\) to subtree \(v\)'' or ``activate \(v\) for its descendants''.  The
difference is only \emph{what you query}: at the same \(\tin(u)\), a subtree
add sees \(x\) when \(u\) is a \emph{descendant} of the source; a root-path
read sees \(x\) when the source is an \emph{ancestor} of \(u\).  Ancestor and
descendant are one relation read from opposite ends -- so one difference BIT
serves both, and this is the whole trick behind CF~383C below.
\end{insightbox}

% ---------------------------------------------------------------
\subsection{Folding a Depth Sign into the Stored Value}
% ---------------------------------------------------------------

\begin{theorem}{Alternating-by-depth becomes a \(\pm1\) multiplier}{c14-parity}
Root the tree; let \(\operatorname{sgn}(v)=+1\) if \(\depth(v)\) is even and
\(-1\) otherwise.  An update ``add \(x\) at node \(r\), and let it propagate
with a sign flip at every edge'' gives node \(u\in\subtree(r)\) the
contribution
\[
  x\cdot(-1)^{\depth(u)-\depth(r)} \;=\; x\cdot\operatorname{sgn}(u)\cdot\operatorname{sgn}(r).
\]
Hence if we store, for each update, the single number \(\Delta=x\cdot\operatorname{sgn}(r)\)
as a \emph{plain} subtree-range add on \([\tin(r),\tout(r)]\), then reading a
node is \(\operatorname{sgn}(u)\cdot(\text{point read at }\tin(u))\), which
reproduces the alternating contribution exactly.  The sign pattern costs
nothing but one \(\pm1\) multiply on each end.
\end{theorem}
\begin{proofbox}
Since \((-1)^{a-b}=(-1)^{a+b}=(-1)^a(-1)^b\) for integers, the propagated
sign of update \(r\) at node \(u\) is \(\operatorname{sgn}(u)\operatorname{sgn}(r)\).
The range add stores \(\Delta=x\operatorname{sgn}(r)\) so that \emph{every}
slot in \([\tin(r),\tout(r)]\) (i.e.\ every descendant of \(r\)) accumulates
\(\Delta\); a point read at \(u\) is the sum \(\sum_{r\,\text{anc. of}\,u}
x_r\operatorname{sgn}(r)\) over exactly the updates whose subtree contains
\(u\).  Multiplying by \(\operatorname{sgn}(u)\) gives
\(\sum_r x_r\operatorname{sgn}(u)\operatorname{sgn}(r)=\sum_r x_r(-1)^{\depth(u)-\depth(r)}\),
the definition of the propagating process.  Nodes outside \(\subtree(r)\) never
see \(\Delta\) because their \(\tin\) is outside its range.
\end{proofbox}

\noindent\textbf{In plain words.}  Do not simulate the up-down-up propagation.
Precompute each node's \(\operatorname{sgn}\) once in the same DFS that stamps
\(\tin/\tout\), then treat ``add to a propagating subtree'' as \emph{one}
difference-BIT range add of a signed amount.  \(\thmref{c14-parity}\) is what
turns CF~383C from a scary statement into \(\coderef{c14-subtree}\).

\begin{tracebox}[Parity store on \(T\): ``add 4 at node 3'' propagates]
\(\depth(3)=1\) so \(\operatorname{sgn}(3)=-1\); \(\Delta=4\cdot(-1)=-4\) added
to range \([\tin_3,\tout_3]=[6,10]\).  Read \(u\) as
\(\operatorname{sgn}(u)\cdot(\text{point at }\tin(u))\), point over the range
\(=-4\) for all of \(\subtree(3)=\{3,6,7,9,10\}\):
\[
\renewcommand{\arraystretch}{1.1}
\begin{array}{c|c|c|c|c}
u & \depth & \operatorname{sgn}(u) & \text{point} & \text{contrib }=\operatorname{sgn}(u)\cdot(-4)\\ \hline
3 & 1 & -1 & -4 & +4\ (\text{self})\\
6 & 2 & +1 & -4 & -4\ (\text{child})\\
7 & 2 & +1 & -4 & -4\\
9 & 3 & -1 & -4 & +4\ (\text{grandchild})\\
10& 4 & +1 & -4 & -4
\end{array}
\]
Self \(+4\), children \(-4\), grandchildren \(+4\), great-grandchild \(-4\) --
the alternating propagation of the statement, from a single range add and a
read-time sign (\(\thmref{c14-parity}\)).
\end{tracebox}

% ---------------------------------------------------------------
\subsection{Problem Lab: CF 383C Propagating Tree}
% ---------------------------------------------------------------

\solvedproblem{c14-prop}{Propagating tree}{CF 383C}
  {https://codeforces.com/problemset/problem/383/C}{CF 2000 L5}
  {Euler tour + one difference BIT, depth-parity sign folded into the stored delta}

\begin{problemstatement}[CF 383C Propagating tree]
A tree of \(n\) nodes numbered \(1..n\), root~\(1\); node \(i\) has value
\(a_i\).  Its \emph{propagating} rule: when a value \(\mathit{val}\) is added
to node \(i\), \(-\mathit{val}\) is added to every child of \(i\); and when
\(-\mathit{val}\) reaches a child, \(+\mathit{val}\) reaches that child's
children, and so on down the whole subtree -- the sign flips every level.
Two query types: \(\texttt{1 }x\,\mathit{val}\) performs the propagating
addition starting at \(x\); \(\texttt{2 }x\) prints the current value at \(x\).

\textbf{Input.} \(n,m\le 2\cdot10^5\); then \(a_1..a_n\) (\(1\le a_i\le1000\));
then \(n-1\) edges; then \(m\) queries with \(1\le x\le n\), \(1\le\mathit{val}\le1000\).

\textbf{Output.} For each type~2 query, one line: the value at \(x\).

\textbf{Example.}
\begin{verbatim}
Input:
5 5
1 2 1 1 2
1 2
1 3
2 4
2 5
1 2 3
1 1 2
2 1
2 2
2 4

Output:
3
3
0
\end{verbatim}
\end{problemstatement}

\begin{attemptgate}[Attempt it yourself first -- \lv{5}, give yourself 30 minutes]
``Add to a subtree with a sign that flips by level'' sounds like you must
walk the subtree.  But the tree is static, so the subtree is a flat interval
(\(\thmref{c24-range}\)) and the sign is a fixed \(\pm1\) per node
(\(\thmref{c14-parity}\)).  Can you make the update \emph{one} range addition?
\peek{store \(\mathit{val}\cdot\!\operatorname{sgn}(x)\) on the interval and read with a sign.}
\nudge{a difference BIT over Euler slots handles ``add to subtree, read one
node'' (\(\coderef{c14-subtree}\) shape B).  All that is left is the parity
factor of \(\thmref{c14-parity}\).}
\fullhint{Precompute \(\tin,\tout,\operatorname{sgn}\) in one DFS.  Query 1
\((x,\mathit{val})\): \(\code{subtree\_add}(x,\;\mathit{val}\cdot\operatorname{sgn}(x))\).
Query 2 \((x)\): print \(a_x+\operatorname{sgn}(x)\cdot\code{node\_value}(x)\).
Both \(\bigO{\log n}\); use \code{long long} for the accumulated values
(\(m\cdot1000\) drifts past \code{int}).}
\end{attemptgate}

\solstep{Understanding}
Two facts collapse the problem.  (1) The tree never changes, so ``the subtree
of \(x\)'' is the fixed index interval \([\tin(x),\tout(x)]\) -- Ch.~5's
\(\thmref{c24-range}\), which we use as a black box.  (2) The propagating sign
is not a process to simulate: it is \(+1/-1\) determined by the parity of the
\emph{depth difference}, which \(\thmref{c14-parity}\) factors into
\(\operatorname{sgn}(u)\operatorname{sgn}(x)\).  Once factored, the update is a
range addition of a signed constant and the read multiplies by a per-node
constant -- \(\propref{c14-eulerrange}\) shape B with one extra \(\pm1\).

\solstep{Key Observations}
\begin{observation}
A type~1 query touches only \(\subtree(x)\); in Euler order that is a
\emph{contiguous} range, so the difference BIT's two writes
(\(+\Delta\) at \(\tin(x)\), \(-\Delta\) at \(\tout(x)+1\)) do the whole
update in \(\bigO{\log n}\).
\end{observation}
\begin{observation}
The node value never needs to be materialised: the answer to a type~2 query is
\(a_x\) (frozen) plus the propagated total, which is
\(\operatorname{sgn}(x)\) times a single point read.
\end{observation}
\begin{observation}
\(m\le2\cdot10^5\) additions of \(\mathit{val}\le1000\) can move a value by
\(\pm2\cdot10^8\), and with alternating signs the accumulator can be negative
-- \code{long long} for the BIT, and signed arithmetic throughout.
\end{observation}

\solstep{Think It Through With Me}
\begin{thinkbox}
\deadend{Recolour the subtree level by level on each query}
{A subtree can be \(\Theta(n)\) and there are \(\Theta(n)\) queries --
\(\bigO{n^2}\) worst case at \(n=2\cdot10^5\), far past the 2\,s budget.  The
propagation is a red herring: it is a \emph{formula}, not a \emph{process}.}
\deadend{Keep two segment trees, one per depth parity}
{Works, but you pay a \(4n\) tree twice and lazy tags for a task that is just
``range add, point read''.  That is the difference BIT's signature job
(\secref{c14-range}); a segment tree is overkill and easy to get wrong.}
\aha{One difference BIT plus a \(\pm1\) sign array does everything: store
\(\mathit{val}\cdot\operatorname{sgn}(x)\) on the subtree interval, and read
\(a_x+\operatorname{sgn}(x)\cdot\text{point}(\tin(x))\).}
\end{thinkbox}

\solstep{The Approach Ladder}
\begin{center}\renewcommand{\arraystretch}{1.25}
\begin{tabular}{@{}c p{0.245\linewidth} p{0.29\linewidth} l c@{}}
\toprule
\textbf{\#} & \textbf{Approach} & \textbf{Idea} & \textbf{Time} & \textbf{Verdict}\\
\midrule
1 & \lv{2}\ BFS the subtree & apply the flip edge by edge & \(\bigO{nm}\) & \tle \\
2 & \lv{3}\ two segment trees & even/odd depth, lazy range add & \(\bigO{n\log n}\) & heavy \\
3 & \lv{5}\ \textbf{difference BIT + parity} & one tree, store \(\mathit{val}\cdot\operatorname{sgn}\) (\(\thmref{c14-parity}\)) & \(\bigO{n\log n}\) & \ac \\
\bottomrule
\end{tabular}\end{center}
Rung~3 is \(\coderef{c14-subtree}\)'s shape B plus \(\thmref{c14-parity}\) --
no new data structure, one extra multiply.

\solstep{Solution}
One DFS fills \(\tin,\tout\) (Ch.~5) and, in the same walk, \(\operatorname{sgn}(v)=(-1)^{\depth(v)}\).
A single difference BIT \(D\) over slots \(1..n\).  Query 1\((x,\mathit{val})\):
\(\code{subtree\_add}(x,\mathit{val}\cdot\operatorname{sgn}(x))\).  Query
2\((x)\): print \(a_x+\operatorname{sgn}(x)\cdot\code{node\_value}(x)\).

\solstep{Proof of Correctness}
\begin{proofbox}
\emph{Range coverage:} \(\code{subtree\_add}\) writes \(+\Delta\) at
\(\tin(x)\) and \(-\Delta\) at \(\tout(x)+1\); its point read at \(u\) equals
\(\Delta\) iff \(\tin(x)\le\tin(u)\le\tout(x)\), i.e.\ iff \(u\in\subtree(x)\)
(\(\thmref{c24-range}\)), and \(0\) otherwise -- matching ``propagation only
reaches descendants''.  \emph{Sign:} the stored delta is
\(\Delta=\mathit{val}\operatorname{sgn}(x)\); the read multiplies by
\(\operatorname{sgn}(u)\), contributing
\(\mathit{val}\operatorname{sgn}(u)\operatorname{sgn}(x)=\mathit{val}(-1)^{\depth(u)-\depth(x)}\),
exactly the level-by-level flip after \(\depth(u)-\depth(x)\) edges
(\(\thmref{c14-parity}\)).  \emph{Superposition:} each query~1 is linear and
independent, so a node's value is \(a_x\) plus the sum of its contributions,
which the prefix sum accumulates.  \emph{Bounds:} \(\bigO{\log n}\) per query,
\(\bigO{n\log n}\) total with the sort-free DFS; \code{long long} is sound
since each value stays within \(\pm(1000\cdot2\cdot10^5+1000)\).
\end{proofbox}

\solstep{Implementation}
\begin{lstlisting}[caption={CF 383C: Euler tour + one difference BIT with a depth-parity sign}, label={lst:c14-383-soln}]
#include <bits/stdc++.h>
using namespace std;
const int MX = 200005;
int n, m, tin[MX], tout[MX], sgn[MX], tmr_;
long long a[MX], D[MX];
vector<int> adj[MX];

void dfs(int u, int p, int dep){                 // Ch.5 tour + parity, one pass
    tin[u] = ++tmr_;  sgn[u] = (dep & 1) ? -1 : 1;
    for (int v : adj[u]) if (v != p) dfs(v, u, dep + 1);
    tout[u] = tmr_;                              // inclusive
}
int lowbit(int i){ return i & -i; }
void bitadd(long long t[], int i, long long d){ for(; i<=n; i+=lowbit(i)) t[i]+=d; }
void range_add(int l, int r, long long d){ bitadd(D,l,d); bitadd(D,r+1,-d); } // 18.2
long long point(int i){ long long s=0; for(; i>0; i-=lowbit(i)) s+=D[i]; return s; }

int main(){
    scanf("%d%d", &n, &m);
    for (int i=1;i<=n;i++) scanf("%lld", &a[i]);
    for (int i=1;i<n;i++){ int u,v; scanf("%d%d",&u,&v);
        adj[u].push_back(v); adj[v].push_back(u); }
    dfs(1, 0, 0);
    while (m--){
        int type, x; scanf("%d%d", &type, &x);
        if (type==1){ long long val; scanf("%lld", &val);
            range_add(tin[x], tout[x], val * sgn[x]);    // signed subtree add
        } else {
            printf("%lld\n", a[x] + sgn[x] * point(tin[x]));
        }
    }
}
\end{lstlisting}

\solstep{Trace on the Official Sample}
\begin{tracebox}[CF 383C: tree \(1\!\to\!2,1\!\to\!3,2\!\to\!4,2\!\to\!5\)]
Root~1, so \(\depth=(0,1,1,2,2)\) and \(\operatorname{sgn}=(+1,-1,-1,+1,+1)\)
for nodes \(1..5\).  \(\tin=(1,2,5,3,4)\), \(\tout=(5,4,5,3,4)\) (subtree\(2=[2,4]\)).
\(a=(1,2,1,1,2)\).

\emph{Q1 \texttt{1 2 3}:} \(\Delta=3\cdot\operatorname{sgn}(2)=-3\) added to
\([2,4]\): \(D\mathrel{+}=(-3\) at~2, \(+3\) at~5\()\).
\emph{Q2 \texttt{1 1 2}:} \(\Delta=2\cdot(+1)=+2\) added to \([1,5]\):
\(D\mathrel{+}=(+2\) at~1, \(-2\) at~6 no-op\()\).
Now the difference array (slots \(1..5\)) is \((+2,-3,0,0,+3)\); prefix sums
\(=(2,-1,-1,-1,2)\).

\emph{Q3 \texttt{2 1}:} point\((\tin_1{=}1)=2\); \(a_1+\operatorname{sgn}(1)\cdot2=1+2=\mathbf{3}\).\\
\emph{Q4 \texttt{2 2}:} point\((\tin_2{=}2)=-1\); \(a_2+\operatorname{sgn}(2)(-1)=2+1=\mathbf{3}\).\\
\emph{Q5 \texttt{2 4}:} point\((\tin_4{=}3)=-1\); \(a_4+\operatorname{sgn}(4)(-1)=1-1=\mathbf{0}\).\\
Printed \(3,3,0\) -- the official output \checkmark.  (Node~5 reads \(a_5+1\cdot(-1)=1\),
node~3 reads \(a_3+(-1)\cdot2=-1\), matching the statement's worked note.)
\end{tracebox}

\solstep{Complexity}
\begin{complexitybox}
One DFS (\(\bigO{n}\)) to fill \(\tin/\tout/\operatorname{sgn}\), then \(m\)
queries at \(\bigO{\log n}\) each: total \(\bigO{(n+m)\log n}\) with
\(n,m\le2\cdot10^5\) -- roughly \(4\cdot10^6\) tree touches, easily inside 2\,s.
Memory: two \(O(n)\) \code{long long} arrays, two \(O(n)\) \code{int} arrays,
and the adjacency list \(2(n-1)\) -- about \(20\) MB, well under 256 MB.
\textbf{AC}.
\end{complexitybox}

\solstep{Verification and Edge Cases}
\begin{center}\renewcommand{\arraystretch}{1.2}
\begin{tabular}{@{}p{0.20\linewidth} p{0.30\linewidth} p{0.38\linewidth}@{}}
\toprule
\textbf{Case} & \textbf{Why tricky} & \textbf{Code behaviour}\\
\midrule
update at a leaf & subtree is one slot & \(\tin=\tout\), range add hits one point \\
\(\tout(x)=n\) & write at \(\tout+1=n+1\) & \code{bitadd} guard \(i\le n\) drops it safely \\
root update & whole array flips sign & \(\operatorname{sgn}\) already encodes it, no special case \\
value goes negative & alternating adds & \code{long long} and signed read \(a_x+\operatorname{sgn}\cdot\text{pt}\) \\
deep chain \(n=2\cdot10^5\) & recursion depth & raise the stack or DFS iteratively \\
\bottomrule
\end{tabular}\end{center}

\begin{hardeningbox}
\begin{itemize}[nosep,leftmargin=1.4em]
  \item \textbf{One-indexed everything:} \(\tin\in[1,n]\), BIT loop
        \(i\le n\), and \(\tout\) \emph{inclusive}.  If you port a \(0\)-based
        Ch.~5 tour (its \(\tout\) is half-open), \emph{shift by one} before
        touching this BIT -- mixing the two conventions is the classic silent
        WA on subtree ranges.
  \item \textbf{Recursion depth:} a chain tree sends \code{dfs} \(2\cdot10^5\)
        frames deep.  Use an explicit stack, or the compiler flag
        \texttt{-Wl,-stack,}, or compile with \code{int} recursion limit raised;
        a stack overflow shows as RE, not TLE.
  \item \textbf{\(\operatorname{sgn}\) is \(\pm1\), not the depth:} the bug is
        writing \(\mathit{val}\cdot\depth\) instead of \(\mathit{val}\cdot\operatorname{sgn}\).
        Verify on the root: its own contribution must be \(+\mathit{val}\).
  \item \textbf{Fast I/O:} \(m\le2\cdot10^5\) lines of output -- buffer it;
        \code{printf}/\code{scanf} as shown, or untied \code{cin}/\code{cout}.
\end{itemize}
\end{hardeningbox}

% ---------------------------------------------------------------
\subsection{Common Pitfalls}
% ---------------------------------------------------------------

\begin{warningbox}[Trap 1: mixing Ch.~5's \(0\)-based tour with this \(1\)-based BIT]
Ch.~5 (\chref{c24}) builds \(\tin\in[0,n)\) with a \emph{half-open}
\(\tout=\tin+\mathit{size}\); this chapter's BIT uses \(\tin\in[1,n]\) and an
\emph{inclusive} \(\tout=\tin+\mathit{size}-1\).  Feeding the first into
\(\code{query}(\tout)-\code{query}(\tin-1)\) is off by one on \emph{both}
ends: subtree sums come out one node too big, and slot~0 indexes past the
array.  Pick one convention per program and translate in exactly one place.
\end{warningbox}

\begin{warningbox}[Trap 2: treating a root-path read as a subtree read]
\(\propref{c14-path}\): the point read at \(\tin(u)\) sums updates at
\(u\)'s \emph{ancestors}, while \(\propref{c14-eulerrange}\)-B's subtree add is
seen by \(u\)'s \emph{descendants}.  They share the two writes but answer
opposite questions.  Getting ``number of active ancestors of \(u\)'' with a
subtree-add semantics (instead of the enter/leave activation) returns the
descendant count -- right number, wrong nodes.
\end{warningbox}

\begin{warningbox}[Trap 3: simulating the propagation]
CF~383C's statement reads like a process.  Do not BFS the subtree flipping
signs; \(\thmref{c14-parity}\) already collapsed it to one signed range add.
The tell that you fell into this trap: your solution has a ``level'' counter
inside the update loop and passes samples but TLEs on a chain.
\end{warningbox}

\begin{warningbox}[Trap 4: forgetting the read-time \(\operatorname{sgn}\)]
Both ends of \(\thmref{c14-parity}\) matter: multiply the stored delta by
\(\operatorname{sgn}(x)\) \emph{and} the read by \(\operatorname{sgn}(u)\).
Dropping the read-time factor makes same-parity and opposite-parity nodes
return values of the wrong sign -- samples with a shallow tree (depth 1--2)
often still pass, hiding the bug until a hidden test.
\end{warningbox}

% ---------------------------------------------------------------
\subsection{Practice Problems}
% ---------------------------------------------------------------

\begin{stuckbox}
\begin{itemize}[nosep,leftmargin=1.4em]
  \item Did you get \(\tin,\tout\) from Ch.~5 by label and \emph{not} re-derive
        them, and are they in the same (1-based, inclusive) convention as the BIT?
  \item Is your update a single range add, or are you secretly walking the
        subtree?
  \item For a sign-propagating problem, is \(\operatorname{sgn}\) computed once
        in the DFS and applied at \emph{both} the write and the read?
  \item Is \(\tout+1\) handled when \(\tout=n\) (the guard, not a crash)?
\end{itemize}
\end{stuckbox}

\begin{practiceset}
  \practice{Danil and a Part-time Job}{CF 877E}
    {https://codeforces.com/problemset/problem/877/E}
    {CF 2000 L4}
    {``Invert the colour of the whole subtree of \(v\); read the colour of
     \(u\)''.  Colour is a bit, and flipping a subtree is a range add
     \(\bmod 2\): flatten with Ch.~5 and keep a difference BIT over
     \([\tin(v),\tout(v)]\), reading \(\code{node\_value}(u)\bmod 2\) (the
     \(\oplus\) running parity).  Exactly \(\coderef{c14-subtree}\) shape B,
     with the value taken mod~2 -- no segment tree needed.}
  \practice{Little Girl and Problem on Trees}{CF 276E}
    {https://codeforces.com/problemset/problem/276/E}
    {CF 2100 L5}
    {``Add \(x\) to every node at distance exactly \(w\) from \(v\) within
     \(v\)'s subtree''.  Nodes at a fixed depth occupy a contiguous block in
     each subtree's Euler span, so keep \emph{one BIT per depth} over the tour
     positions and range-add \(x\) on depth \(\depth(v)+w\) inside
     \([\tin(v),\tout(v)]\); a node read is one point query in its depth's BIT.
     The lesson: when the query keys on (subtree, depth), index the Euler
     array by depth groups -- the same \(\tin/\tout\) reduction, refined.}
\end{practiceset}

\paragraph{Back-cites and forwards.}
The \(\tin/\tout\) stamps and ``subtree \(=\) interval'' come wholly from
Ch.~5 (\chref{c24}, \(\thmref{c24-range}\)); we supply only the range engine it
deferred.  CSES~1137 ``Subtree Queries'' is Ch.~5's own solved case with a
static prefix array -- replay it here with \(\coderef{c14-subtree}\) shape A and
updates become free.  CF~343D ``Water Tree'' is the same subtree-add /
root-path-read pair with the URL owned by Ch.~6 (\chref{c06}) -- cited by name,
not re-linked.  Forward: ``add on one axis, read on another, offline'' returns
in \secref{c14-offline}; the k-th-by-lifting predicate the Summit CF~1181D
binary-searches over is this section's structure seen from the value side
(\(\secref{c14-order}\)).

% ---------------------------------------------------------------
\begin{recallbox}[Section 18.4 in 10 seconds]
Ch.~5 flattened the tree; a BIT now owns its ranges.  Node\(\,\to\,\)slot
\(\tin(v)\), subtree\(\,\to\,\)interval \([\tin(v),\tout(v)]\)
(\(\defref{c14-eulerbit}\)).  Point-update + subtree-sum is the plain BIT,
subtree-add + point-read is the difference BIT (\(\propref{c14-eulerrange}\),
\(\coderef{c14-subtree}\)); the same two writes read the other way give
root-path sums (\(\propref{c14-path}\), \(\coderef{c14-path}\)).  A depth-parity
sign folds into the stored delta so a propagating update stays one range add
(\(\thmref{c14-parity}\)) -- that solves CF~383C
(\(\probref{c14-prop}\), \(\coderef{c14-383-soln}\)).  Keep one tour
convention per program.
\selftest{
  \item State, in one line each, the reduction for ``subtree sum'' and for
        ``root-path sum'' on a static tree.
  \item Why can \emph{one} difference BIT answer both, and what differs: the
        writes or the read?
  \item For a query ``add to subtree with a sign flipping by level'', what is
        stored and what is the node read?
  \item Give the exact symptom of mixing a \(0\)-based half-open tour with a
        \(1\)-based inclusive BIT.
  \item Which chapter owns \(\tin/\tout\), and which one owns the range engine?
}
\end{recallbox}
